\section{本讲小结}

球面成像的一阶关系来自当地法线与近轴方向的连接。折射采用 $n_1/u+n_2/v=(n_2-n_1)/R_s$，每次先按传播方向判断物、像和球心的位置，再赋予距离符号。反向传播需要交换介质并重定曲率号；平面折射是曲率趋零的极限，沿法线附近观察时才具有统一视深。

多界面问题先保留中间像的实际坐标，再对下一面重新量物距。会聚点尚未到达下一面时，它是下一面的虚物；前一级得到负像距，也不能跳过后续液面而直接当作空气中的最终像。横向放大率、轴向位移放大率与像速分别对应不同几何量，应按各自的定义计算。

镜面反射使用实物、实像为正的约定，实像位于镜前。焦点处的平行状态改变像的虚实，二倍焦距处改变放大与缩小，两个分界应区分。球面镜的 $|f|=R/2$ 只在近轴极限成立；增加孔径后应以实际反射点和法线追迹，才能看见纵向球差以及局部焦点的移动。
